\clearpage\section{Аналитическая часть}
\subsection{Постановка задачи}
Дано полутоновое изображение шириной $W$ и высотой $H$, $W,H>0$. Пиксель --- элемент изображения; его яркость задаётся целым числом от 0 до 255. Число пикселей $N=W\cdot H$. Матрица хранится построчно в массиве $I=(I_0,\ldots,I_{N-1})$.

Требуется выбрать порог $T$ методом Оцу и получить бинарную маску $B$, то есть изображение с двумя значениями яркости:
\begin{equation}\label{eq:mask}
B_i=\begin{cases}0,&I_i\le T,\\255,&I_i>T,\end{cases}\qquad 0\le i<N.
\end{equation}
Нулём обозначается класс тёмных пикселей, значением 255 --- класс светлых. Связность областей и принадлежность конкретному анатомическому объекту не устанавливаются. Под чёрно-белым снимком понимается полутоновый исходник; уже бинарный вход также допустим.

\subsection{Метод Оцу}
Гистограмма яркостей $h[v]$ --- число пикселей со значением $v$. Метод Оцу основан на выборе порога, обеспечивающего наибольшее различие двух классов по яркости~\cite{otsu}. Для вычислений вводятся накопленные количества и суммы:
\begin{equation}\label{eq:sums}
\begin{aligned}
n_0(t)&=\sum_{v=0}^{t}h[v], & n_1(t)&=N-n_0(t),\\
s_0(t)&=\sum_{v=0}^{t}v\cdot h[v], & s&=\sum_{v=0}^{255}v\cdot h[v].
\end{aligned}
\end{equation}
Здесь $n_0,n_1$ --- размеры классов, $s_0$ --- сумма яркостей тёмного класса, $s$ --- сумма яркостей всего изображения. Средние яркости классов равны $\mu_0=s_0/n_0$ и $\mu_1=(s-s_0)/n_1$, если оба класса непусты. Межклассовая дисперсия определяется как
\begin{equation}
\sigma_B^2(t)=\frac{n_0(t)}{N}\cdot\frac{n_1(t)}{N}\cdot\bigl(\mu_0(t)-\mu_1(t)\bigr)^2.
\end{equation}
Для поиска максимума используется эквивалентный критерий без постоянного множителя $N^{-2}$:
\begin{equation}\label{eq:criterion}
J(t)=\frac{\bigl(N\cdot s_0(t)-n_0(t)\cdot s\bigr)^2}{n_0(t)\cdot n_1(t)},\qquad n_0(t)\cdot n_1(t)>0.
\end{equation}
Пороги перебираются от 0 до 254. Накопленные суммы обновляются за постоянное число действий на порог. При одинаковом значении критерия сохраняется меньший порог. Это принятое в работе правило разрешения равенства. Если изображение однотонное, непустое разбиение невозможно: порог принимается равным единственной яркости, маска получается нулевой.

\subsection{Возможности распараллеливания}
Построение гистограммы и применение готового порога требуют обхода всех пикселей. При непосредственном изменении общей гистограммы потребовалась бы синхронизация каждого приращения. Вместо этого каждый рабочий поток формирует собственные 256 счётчиков; после завершения этой фазы диспетчер суммирует их.

Порог зависит от всех локальных гистограмм. Поэтому бинаризацию нельзя начинать до окончания их объединения и поиска порога. После публикации порога каждый рабочий поток обрабатывает назначенный ему диапазон пикселей. Диапазоны разных потоков не перекрываются, поэтому два потока не записывают результат в один и тот же пиксель маски. Используется параллелизм по данным со статическим распределением пикселей.

Метод Оцу требует только гистограмму и не нуждается в заданном вручную пороге. Однако при неоднородном освещении и перекрывающихся распределениях яркости одной границы может быть недостаточно. В работе сравниваются реализации выбранной бинарной сегментации; точность выделения органов по клинической разметке не оценивается.
